%=============================================================
% Paper 5: Lazear & Spletzer (2012), AER P&P
%=============================================================
\chapter{Lazear \& Spletzer (2012)}

\begin{paperbox}{Hiring, Churn, and the Business Cycle}
\textbf{Authors:} Edward P.\ Lazear \& James R.\ Spletzer\\[3pt]
\textbf{Journal:} \textit{American Economic Review: Papers \& Proceedings}
102(3): 575--579\\[3pt]
\textbf{Year:} 2012 \quad\textbf{Field:} Labour economics; Macroeconomics\\[3pt]
\textbf{Classification:} Empirical (descriptive accounting decomposition)
\end{paperbox}

%------------------------------------------------------------
\section{Research Question}

How does employment churn---the simultaneous hiring and separation within a
business that offsets each other without changing net employment---vary over
the business cycle, and what role did changes in churn play in the collapse
of hiring during the 2007--09 Great Recession?

%------------------------------------------------------------
\section{Motivation}

Standard analyses of the Great Recession focused on net employment (the
dramatic fall in payrolls) and on job creation and destruction. But these
net measures obscure a potentially important channel: the cessation of
\emph{churn}. Even businesses that are not changing their employment level
engage in substantial ongoing hiring and separation activity---workers quit,
are hired as replacements, transfer, or retire. This churn is intrinsically
valuable: it enables workers to move to more productive jobs and enables firms
to replace mismatched workers. If churn collapses during recessions, labour
market efficiency losses may substantially exceed what net employment changes
alone reveal. Lazear and Spletzer construct the first explicit churn
decomposition from JOLTS microdata to quantify this channel.

%------------------------------------------------------------
\section{Main Contribution}

The paper introduces an explicit accounting decomposition of total hires and
separations into four components---growth hires ($H^G_E$), replacement hires
(churn in expanding firms), replacement separations (churn in contracting
firms), and employment-decreasing separations---and uses JOLTS microdata to
estimate each component quarterly. Key findings: churn is (i) procyclical
($\rho(\text{churn},u) = -0.96$); (ii) quantitatively enormous (65\% of all
hires in normal times were churn); and (iii) the dominant component of the
hiring collapse during the 2007--09 recession (\textbf{79\%} of the decline
in hires was due to reduced churn, not reduced job creation).

%------------------------------------------------------------
\section{Relationship to the Literature}

The paper fits into the wider gross flows literature initiated by DH1992 but
focuses on within-firm turnover (churn) rather than across-firm flows (job
creation/destruction). It is related to Anderson--Meyer (1994),
Burgess--Lane--Stevens (2000), and Davis--Faberman--Haltiwanger (2006), who
had estimated that 42--70\% of hiring was churn using various data sources.
The paper's contribution is to use JOLTS \emph{microdata}---enabling the
within-establishment decomposition---and to examine the business-cycle
variation directly. Churn is distinct from, but related to, the ``excess
reallocation'' concept in DH1992.

%------------------------------------------------------------
\section{Theoretical Framework and Economic Mechanism}

\subsection*{Accounting Framework}

For an establishment in period $t$, define hires $H$ and separations $S$ in
business of type $j\in\{E=\text{expanding},\;C=\text{contracting},\;
Z=\text{zero-growth}\}$. Total hires and separations decompose as:
\begin{align}
  H &= H^G_E + \text{CH}, \label{eq:hires_decomp} \\
  S &= S^D_C + \text{CH}, \label{eq:seps_decomp}
\end{align}
where $\text{CH}=\text{CH}_E+\text{CH}_C+\text{CH}_Z$ is total churn:
\begin{align}
  \text{CH}_E &= H^R_E = S_E
    &&\text{(replacement hires in expanding firms)} \\
  \text{CH}_C &= H_C = S^R_C
    &&\text{(replacement separations in contracting firms)} \\
  \text{CH}_Z &= H_Z = S_Z
    &&\text{(all hires/seps in zero-growth firms)}
\end{align}
Net employment change follows directly:
\begin{equation}
  H - S = H^G_E - S^D_C.
  \label{eq:net_change}
\end{equation}
\textbf{Crucially, total churn CH cancels exactly from the net change
equation.} Churn has no direct effect on net employment---but it determines
the \emph{level} of total hiring activity.

\subsection*{Two Cyclical Mechanisms}

\begin{enumerate}
  \item \textbf{Demand side:} In downturns, businesses decrease their hiring
    and separations that would have been matched by replacement hires remain
    unfilled $\Rightarrow$ churn falls as unfilled separations become
    employment-decreasing separations.
  \item \textbf{Supply side:} Workers become reluctant to quit in recessions
    (fear of unemployment) $\Rightarrow$ quit rate falls $\Rightarrow$ fewer
    replacement hires needed $\Rightarrow$ churn falls through a different
    channel.
\end{enumerate}

%------------------------------------------------------------
\section{Data}

\begin{itemize}
  \item \textbf{Dataset:} JOLTS (Job Openings and Labour Turnover Survey)
    microdata, U.S.\ Bureau of Labour Statistics.
  \item \textbf{Coverage:} Random sample of $\approx$16,000 nonfarm business
    establishments; $\approx$10,500 provide data regularly.
  \item \textbf{Period:} December 2000 -- June 2011; estimated at quarterly
    frequency.
  \item \textbf{Key feature:} JOLTS microdata allow classification of each
    establishment by its net employment change each quarter (expanding,
    contracting, zero-growth), enabling the within-establishment decomposition.
  \item \textbf{Limitation:} Excludes imputed hires and separations for
    unobserved births and deaths; quarterly estimates ($\approx$12\,m hires in
    mid-2000s) are below published monthly statistics ($\approx$15\,m) partly
    for this reason.
\end{itemize}

%------------------------------------------------------------
\section{Empirical Strategy}

Pure descriptive decomposition. The strategy:
\begin{enumerate}
  \item Classify each JOLTS establishment in each quarter as expanding
    ($\Delta E>0$), contracting ($\Delta E<0$), or zero-growth ($\Delta E=0$).
  \item Within each category, compute hires and separations from JOLTS
    microdata.
  \item Apply the churn accounting identities
    (\ref{eq:hires_decomp})--(\ref{eq:net_change}) to decompose total hires
    and separations.
  \item Compute seasonally adjusted time series for each component (Figure~1).
  \item Formally decompose the change in churn from 2007:Q4 to 2009:Q1 into:
    (a) change in the number of establishments with churn; (b) change in
    average size of churn conditional on having churn (Table~1).
\end{enumerate}

%------------------------------------------------------------
\section{Identification Strategy}

\textbf{No causal identification.} This is a descriptive accounting exercise.
The decomposition is an identity, not a regression. The paper makes no claim
that churn \emph{causes} employment changes or that the recession \emph{caused}
the decline in churn---it documents the accounting fact that 79\% of the
decline in hires during 2007--09 was attributable to a decline in churn. The
welfare implications (cost $\approx0.4$ pp of GDP annually) rely on auxiliary
assumptions about the productivity value of worker reallocation via churn, not
on causal identification.

%------------------------------------------------------------
\section{Main Results}

\subsection*{Stylised Facts about Churn}

\begin{center}
\begin{tabular}{lc}
\toprule
\textbf{Statistic} & \textbf{Value} \\
\midrule
Share of hires that is churn (mid-2000s) & 65\% \\
$\rho(\text{churn},\,u)$ & $-0.96$ \\
$\rho(\text{churn},\,\text{NET})$ & $+0.43$ (2001--2011:Q2) \\
$\rho(\text{churn},\,\text{NET})$ & $+0.79$ (2001--2009:Q2) \\
\bottomrule
\end{tabular}
\end{center}

\subsection*{Great Recession Decomposition}

\begin{center}
\begin{tabular}{lcc}
\toprule
\textbf{Component} & \textbf{2007:Q4} & \textbf{2009:Q2} \\
\midrule
Total churn (millions) & 8.3 & 5.3 \\
Total hires (millions) & $\approx$12 & $\approx$8.2 \\
\midrule
\multicolumn{3}{l}{\textit{Decomposition of the 3.0\,m decline in churn}} \\
Fewer establishments with churn & 1.9\,m $\to$ 1.6\,m & 52\% \\
Smaller average churn size & 4.44 $\to$ 3.78 per estab. & 48\% \\
\midrule
Share of hiring decline due to churn & \multicolumn{2}{c}{\textbf{79\%}} \\
Share due to lower job creation ($H^G_E$) & \multicolumn{2}{c}{$\approx25\%$} \\
\bottomrule
\multicolumn{3}{l}{\small Source: JOLTS microdata, seasonally adjusted.}
\end{tabular}
\end{center}

%------------------------------------------------------------
\section{Economic Magnitude and Interpretation}

The finding that 79\% of the Great Recession's hiring collapse was churn-driven
is the headline result and represents a major reframing of what happened in
the 2007--09 recession. Standard analyses focused on job creation collapsing
and job destruction rising. Lazear--Spletzer show that these forces, while
real, were quantitatively swamped by the simultaneous collapse of the
replacement-hire/quit mechanism. The economy became \emph{frozen} rather than
simply contracting in the traditional sense.

The welfare cost estimate:
\begin{itemize}
  \item Reduced churn $\approx 0.4$ percentage points of GDP \emph{annually}
    over the 3.5 years following the recession onset.
  \item This estimate uses the auxiliary assumption that churn-driven worker
    reallocation raises productivity by observed average productivity differences
    between gaining and losing firms.
\end{itemize}

%------------------------------------------------------------
\section{Mechanisms}

Two mechanisms are empirically consistent but the paper cannot distinguish
between them:

\begin{enumerate}
  \item \textbf{Supply side (worker reluctance):} Workers become reluctant to
    quit in recessions (fear of being unable to find new employment), so the
    quit rate falls. This reduces replacement hiring even at firms that would
    gladly replace a quitter.

  \item \textbf{Demand side (employer frugality):} Employers, when a worker
    separates, choose not to replace them---tolerating temporary staffing
    gaps. This mechanism is the paper's preferred interpretation: ``separations
    that would have been matched by hiring during good times remain unfilled.''
    The causal chain runs: aggregate shock $\to$ employer decision not to
    replace $\to$ reduced churn.
\end{enumerate}

The 52/48 extensive/intensive margin split of the churn decline is consistent
with both mechanisms.

%------------------------------------------------------------
\section{Robustness Checks}

\begin{itemize}
  \item Validation of JOLTS microdata estimates against published JOLTS
    statistics (hires and separations track published series).
  \item Validation of job creation/destruction estimates against BLS Business
    Employment Dynamics (BED) data.
  \item Consistency of the 65\% churn share with prior estimates (Anderson--Meyer
    1994: 69\%; Burgess--Lane--Stevens 2000: $\approx$70\%; Davis et al.\ 2006:
    42--72\%).
\end{itemize}

Note that exclusion of births/deaths imputation is a consistent understatement
and does not affect the proportional decomposition.

%------------------------------------------------------------
\section{Limitations}

\begin{enumerate}
  \item \textbf{JOLTS sample:} $\approx$10,500 establishments, disproportionately
    large businesses. Small businesses, which tend to have higher churn rates,
    may be underrepresented.
  \item \textbf{Short time series:} JOLTS microdata begin in December 2000;
    only one full business cycle is available, with the 2001 recession only
    partially captured.
  \item \textbf{No worker-level data:} Churn is measured at the establishment
    level. It is impossible to identify whether replacement hires are the same
    or different workers from those who separated---which matters for
    productivity analysis.
  \item \textbf{Nonfarm private sector only:} Government employment is
    excluded; education and health services differ substantially in quit
    behaviour.
  \item \textbf{No causal identification:} The paper documents the accounting
    fact that churn declined. Whether this represents an efficiency loss or
    a rational adjustment to changed incentives is not answered.
  \item \textbf{P\&P format:} As a 5-page Papers \& Proceedings paper, many
    analytical details are necessarily compressed. This is a short report of
    findings rather than a full research paper.
\end{enumerate}

%------------------------------------------------------------
\section{Policy Implications}

\begin{itemize}
  \item If 79\% of the Great Recession's hiring collapse was churn-related,
    then policies designed to stimulate job creation (investment tax credits,
    hiring subsidies) would have addressed only $\approx$25\% of the hiring
    shortfall.
  \item Policies that encourage worker mobility (portable benefits, reductions
    in non-wage costs of switching jobs) would target the mechanism directly.
  \item The $\approx$0.4\% of GDP welfare cost from reduced churn provides a
    rough benchmark for evaluating policies that restore labour market
    fluidity.
  \item Aggregate demand stimulus might restore churn partly by restoring
    worker confidence to quit---a labour supply channel distinct from the job
    creation/destruction channel.
\end{itemize}

%------------------------------------------------------------
\section{Overall Assessment}

This is an insightful, well-executed descriptive paper that reframes a key
fact about the Great Recession. The accounting decomposition is clean and the
empirical finding is striking. Its limitations as a short P\&P paper are
acknowledged: it establishes a stylised fact rather than a causal mechanism,
and the welfare calculation rests on auxiliary assumptions that are not formally
derived. The paper raises more questions than it answers---the mechanisms
behind the churn collapse, the welfare implications, the heterogeneity across
industries and workers---which is entirely appropriate for a conference
proceedings contribution that motivates a more complete follow-up.

\textbf{Quality: Good for its format. Concise, clear, empirically novel,
and policy-relevant.}

%------------------------------------------------------------
\section{Relevance to My Research}

Very relevant for RDC work on Canadian labour dynamics. The JOLTS-based
decomposition could be replicated using Canadian administrative data (LEED or
the Longitudinal Worker File), which allows similar classification of
establishments by their employment direction. The hypothesis that churn is
procyclical and was severely disrupted during the 2008--09 recession and the
COVID-19 pandemic in Canada could be tested directly. Of particular interest
would be whether the sectoral composition of churn collapse differs in Canada:
given the resource sector's sensitivity to commodity cycles, churn dynamics in
Alberta and Saskatchewan would likely differ sharply from Ontario and
Qu\'{e}bec during commodity busts.

%------------------------------------------------------------
\section{Potential Research Extensions}

\begin{itemize}
  \item Replicate the Lazear--Spletzer decomposition for Canada using
    LWF/LEED data for 2000--2024, covering the 2008--09 recession and the
    COVID-19 shock.
  \item Disaggregate churn by industry, province, worker age, and occupation
    to identify which segments of the Canadian labour market experienced the
    largest efficiency losses from churn collapse.
  \item Estimate a structural model of the quit/hire decision to formally
    identify whether the churn collapse in Canada was supply-driven (workers
    reluctant to quit) or demand-driven (employers declining to replace
    quitters), using variation in provincial unemployment insurance generosity
    as an instrument for worker quit behaviour.
\end{itemize}

%------------------------------------------------------------
\section{Research Gaps}

\begin{itemize}
  \item The paper uses establishment-level data but cannot identify whether
    churn involves high-to-low or low-to-high productivity transitions. A key
    question is whether \emph{efficient} churn collapsed during the recession
    or whether the decline was concentrated in inefficient (mismatch-correcting)
    turnover.
  \item The cost estimate (0.4\% of GDP/year) is rough and does not account
    for the possibility that some churn during booms was excessive
    (overhiring/overquitting driven by optimism), so not all lost churn
    represents an efficiency loss.
\end{itemize}

%------------------------------------------------------------
\section{Methodological Lessons}

\begin{enumerate}
  \item A simple accounting decomposition can reveal economically important
    facts invisible in aggregate statistics---the 79\% finding emerges directly
    from the churn identity applied to microdata and would be invisible in
    aggregate employment or JOLTS-level statistics.
  \item Classifying establishments by employment \emph{direction}
    (expanding/contracting/zero-growth) within each period, rather than by
    long-run growth status, is essential for implementing the churn
    decomposition correctly.
  \item Validating microdata-based estimates against published aggregate
    statistics (BED, published JOLTS) provides reassurance that microdata
    classifications are consistent with external benchmarks.
\end{enumerate}

%------------------------------------------------------------
\section*{Five-Sentence Summary}

Lazear and Spletzer (2012) use JOLTS establishment microdata to decompose
total hires and separations into four components: growth hires (job creation),
employment-decreasing separations (job destruction), and churn (replacement
hires and separations that offset each other within a firm without changing
net employment). They show that churn accounts for approximately 65\% of all
hiring in normal times and is strongly procyclical, with a correlation of
$-0.96$ with the unemployment rate over 2001--2011. During the 2007--09 Great
Recession, churn fell by 3.0 million workers per quarter (36\% decline), and
this churn collapse accounts for 79\% of the total decline in hiring---far
larger than the decline in job creation. The decline in churn was evenly split
between fewer establishments experiencing any churn (52\%) and smaller average
churn among those that did (48\%), consistent with a combination of worker
reluctance to quit and employer decisions to leave separations unfilled. The
welfare cost of this reduced churn---lost productivity from foregone efficient
worker reallocation---is estimated at roughly 0.4 percentage points of GDP
annually over the post-recession period, highlighting that the efficiency costs
of recessions substantially exceed those revealed by net employment changes
alone.

\vspace{6pt}
\begin{quickref}
\begin{tabularx}{\linewidth}{@{}lX@{}}
\textbf{Key contribution} &
  First explicit churn decomposition from JOLTS microdata; documents that
  79\% of the Great Recession's hiring collapse was churn-driven rather than
  job-creation-driven---fundamentally reframing the labour market contraction. \\[4pt]
\textbf{Key weakness} &
  Short P\&P format; no causal identification; cannot distinguish supply-side
  (reluctant quitters) from demand-side (non-replacing employers) mechanisms;
  no worker-level data to identify productivity effects of churn. \\[4pt]
\textbf{Most interesting gap} &
  Whether the churn collapse was efficient (eliminating mismatch-driven
  excessive turnover) or inefficient (preventing productive reallocation)---the
  answer determines whether recession-period policy should try to restore churn
  or accept it as rational adjustment. \\[4pt]
\textbf{Publishable extension} &
  A full-length paper using Canadian LWF/LEED data to replicate the churn
  decomposition over 2000--2024, disaggregating by province, industry, age,
  and occupation, and formally testing the supply vs.\ demand mechanism using
  provincial unemployment insurance generosity as an instrument for worker
  quit behaviour.
\end{tabularx}
\end{quickref}
